\documentclass[11pt,a4paper,twoside]{article}

\usepackage[a4paper,top=2.25cm,bottom=2.25cm,inner=2.35cm,outer=1.85cm,headheight=15pt]{geometry}
\usepackage{fontspec}
\usepackage{xeCJK}
\setmainfont{Times New Roman}
\IfFontExistsTF{Microsoft YaHei}
  {\setCJKmainfont{Microsoft YaHei}\setCJKsansfont{Microsoft YaHei}}
  {\setCJKmainfont{SimSun}\setCJKsansfont{SimHei}}
\usepackage{amsmath,amssymb}
\usepackage{xcolor}
\usepackage{enumitem}
\usepackage{titlesec}
\usepackage[most]{tcolorbox}
\usepackage{fancyhdr}
\usepackage{lastpage}
\usepackage{hyperref}
\usepackage{bookmark}
\usepackage{needspace}

\definecolor{accent}{RGB}{30,70,110}
\definecolor{enback}{RGB}{246,249,252}
\definecolor{zhback}{RGB}{250,250,250}
\definecolor{rulegray}{RGB}{120,120,120}

\hypersetup{
  unicode=true,
  colorlinks=true,
  linkcolor=accent,
  pdftitle={Bilingual Reviewer Comments / 审稿意见中英对照},
  pdfsubject={NETJOURNAL-D-26-00829 reviewer comments only}
}

\pagestyle{fancy}
\fancyhf{}
\fancyhead[LE,RO]{\small\color{rulegray}审稿意见中英对照}
\fancyhead[RE,LO]{\small\color{rulegray}NETJOURNAL-D-26-00829}
\fancyfoot[C]{\small\thepage\,/\,\pageref*{LastPage}}
\renewcommand{\headrulewidth}{0.35pt}
\renewcommand{\footrulewidth}{0pt}

\titleformat{\section}{\Large\bfseries\color{accent}}{}{0pt}{}
\titleformat{\subsection}{\large\bfseries}{}{0pt}{}
\titlespacing*{\section}{0pt}{1.2em}{0.55em}
\titlespacing*{\subsection}{0pt}{1.05em}{0.45em}
\setlength{\parindent}{0pt}
\setlength{\parskip}{0.4em}
\setlist[enumerate]{leftmargin=2.2em,itemsep=0.25em,topsep=0.25em}
\setcounter{tocdepth}{2}
\setlength{\emergencystretch}{2em}

\newtcolorbox{englishcomment}{
  enhanced,breakable,
  colback=enback,colframe=accent,
  boxrule=0.55pt,arc=1.5mm,
  left=2.5mm,right=2.5mm,top=2mm,bottom=2mm,
  before skip=0.25em,after skip=0.55em,
  fonttitle=\bfseries,
  title={English original}
}
\newtcolorbox{chinesecomment}{
  enhanced,breakable,
  colback=zhback,colframe=rulegray,
  boxrule=0.55pt,arc=1.5mm,
  left=2.5mm,right=2.5mm,top=2mm,bottom=2mm,
  before skip=0.25em,after skip=0.9em,
  fonttitle=\bfseries,
  title={中文译文}
}

\begin{document}

\begin{center}
  {\LARGE\bfseries Bilingual Reviewer Comments}\\[0.35em]
  {\LARGE\bfseries 审稿意见中英对照}\\[0.9em]
  \begin{tabular}{@{}rl@{}}
    \textbf{Manuscript / 稿件编号：} & NETJOURNAL-D-26-00829 \\
    \textbf{Journal / 期刊：} & Nuclear Engineering and Technology \\
  \end{tabular}
\end{center}

\vspace{0.55em}
\begin{tcolorbox}[colback=white,colframe=rulegray,boxrule=0.45pt,arc=1mm,left=2mm,right=2mm,top=1.5mm,bottom=1.5mm]
\small 本文件仅汇编审稿人意见；每条先列英文原文，再列中文译文，不含作者回复与修改说明。\\
\textit{This document contains reviewer comments only. Each English original is followed by its Chinese translation; author responses and revision notes are omitted.}
\end{tcolorbox}

\tableofcontents
\clearpage
\Needspace{14\baselineskip}
\section{Reviewer \#1 / 审稿人 \#1}
\Needspace{8\baselineskip}
\subsection{R1.P: Reviewer overview / 审稿人总体评价}
\begin{englishcomment}
This paper proposes a deep learning framework for solving neutron diffusion equations in heterogeneous media. The primary challenge addressed is the presence of discontinuous diffusion coefficients and macroscopic cross sections at material interfaces. These discontinuities introduce flux-gradient jumps and interface singularities that may adversely affect the accuracy and convergence of neural-network-based solvers. To address this issue, the authors formulate the neutron diffusion equation using a Variable-Domain Integral Weak Solution (VIWS) approach and employ a single-network representation of the global solution. The proposed methodology aims to solve heterogeneous neutron diffusion problems without requiring explicit treatment of internal material interfaces. The topic is relevant and potentially valuable for the application of deep learning methods in reactor physics. However, the manuscript lacks sufficient methodological and implementation details to allow readers to fully understand, reproduce, and evaluate the proposed approach. The following comments should be addressed.
\end{englishcomment}
\begin{chinesecomment}
本文提出了一种用于求解非均匀介质中中子扩散方程的深度学习框架。本文所处理的主要挑战是材料界面处存在不连续的扩散系数和宏观截面。这些不连续性会引起通量梯度跳跃和界面奇异性，可能对基于神经网络的求解器的精度与收敛性产生不利影响。为解决这一问题，作者采用变域积分弱解（Variable-Domain Integral Weak Solution, VIWS）方法建立中子扩散方程，并使用单网络表示全局解。该方法旨在无需显式处理内部材料界面的情况下求解非均匀中子扩散问题。该主题与反应堆物理中深度学习方法的应用相关，并具有潜在价值。然而，稿件缺少足够的方法与实现细节，读者尚难以充分理解、复现和评价所提出的方法。以下意见需要处理。
\end{chinesecomment}

\Needspace{8\baselineskip}
\subsection{R1.1: Comment 1 / 意见 1}
\begin{englishcomment}
The abstract does not clearly identify the comparison baseline. It would be beneficial to explicitly state whether the proposed method is compared against conventional PINNs, finite element methods (FEM), finite difference methods (FDM), or other established neutron diffusion solvers. In addition, the abstract does not provide any indication of computational cost, training time, memory requirements, or scalability to large-scale reactor calculations.
\end{englishcomment}
\begin{chinesecomment}
摘要没有清楚说明比较基线。建议明确指出所提出的方法是与传统 PINN、有限元方法（FEM）、有限差分方法（FDM），还是其他成熟的中子扩散求解器进行比较。此外，摘要没有说明计算成本、训练时间、内存需求，或向大规模反应堆计算扩展的能力。
\end{chinesecomment}

\Needspace{8\baselineskip}
\subsection{R1.2: Comment 2 / 意见 2}
\begin{englishcomment}
The manuscript should provide a more detailed discussion regarding the existence and uniqueness conditions associated with the antiderivative formulation employed in the proposed methodology.
\end{englishcomment}
\begin{chinesecomment}
稿件应更详细地讨论所提出方法中原函数形式所涉及的存在性与唯一性条件。
\end{chinesecomment}

\Needspace{8\baselineskip}
\subsection{R1.A: Unnumbered architecture comment / 原决定信中未编号的网络结构意见}
\begin{englishcomment}
The neural-network architecture is insufficiently described. If separate antiderivative and neutron-flux networks are employed, their structures should be explicitly specified. At a minimum, the following information should be provided:
\begin{enumerate}[label=\alph*.,leftmargin=2em]
\item Network structure
\item Number of hidden layers
\item Number of neurons per layer
\item Activation functions
\item Input/output variables
\item Training strategy
\end{enumerate}
\end{englishcomment}
\begin{chinesecomment}
稿件对神经网络结构的描述不够充分。如果分别使用原函数网络和中子通量网络，应明确说明它们的结构。至少应提供以下信息：
\begin{enumerate}[label=\alph*.,leftmargin=2em]
\item 网络结构；
\item 隐藏层数量；
\item 每层神经元数量；
\item 激活函数；
\item 输入/输出变量；
\item 训练策略。
\end{enumerate}
\end{chinesecomment}

\Needspace{8\baselineskip}
\subsection{R1.3: Comment 3 / 意见 3}
\begin{englishcomment}
The optimization procedure is not adequately described. The manuscript should clearly state the trainable parameters, optimization algorithm(s), learning-rate strategy, and any additional optimization settings used during training.
\end{englishcomment}
\begin{chinesecomment}
优化过程的描述不够充分。稿件应清楚说明可训练参数、优化算法、学习率策略，以及训练过程中采用的其他优化设置。
\end{chinesecomment}

\Needspace{8\baselineskip}
\subsection{R1.4: Comment 4 / 意见 4}
\begin{englishcomment}
The manuscript states that training terminates when the loss reaches a prescribed accuracy requirement or satisfies a stopping criterion. However, neither the accuracy requirement nor the stopping criterion is explicitly defined. Since References 38 and 39 are cited for these details, and one of these references is not readily accessible to an international audience, the relevant information should be included directly in the manuscript.
\end{englishcomment}
\begin{chinesecomment}
稿件称，当损失达到预设精度要求或满足停止准则时终止训练。然而，文中既没有明确定义该精度要求，也没有明确定义停止准则。稿件引用参考文献 38 和 39 说明这些细节，但其中一篇文献不便于国际读者获取，因此相关信息应直接写入稿件。
\end{chinesecomment}

\Needspace{8\baselineskip}
\subsection{R1.5: Comment 5 / 意见 5}
\begin{englishcomment}
The term L-BFGS appears in Section 4 without prior introduction. The authors should provide the full name and briefly explain its role within the optimization procedure.
\end{englishcomment}
\begin{chinesecomment}
术语 L-BFGS 在第 4 节出现之前没有介绍。作者应给出其全称，并简要说明它在优化过程中的作用。
\end{chinesecomment}

\Needspace{8\baselineskip}
\subsection{R1.6: Comment 6 / 意见 6}
\begin{englishcomment}
The manuscript states that the ``default network is an 8-layer FCNN.'' The meaning of ``default'' should be clarified. Is this architecture used throughout all numerical experiments, or only when otherwise unspecified?
\end{englishcomment}
\begin{chinesecomment}
稿件称“默认网络为 8 层 FCNN”。应澄清“默认”的含义：该结构是否用于全部数值实验，还是仅在没有另行说明时使用？
\end{chinesecomment}

\Needspace{8\baselineskip}
\subsection{R1.7: Comment 7 / 意见 7}
\begin{englishcomment}
The manuscript should explicitly state the number and physical meaning of the input and output variables used by each neural network.
\end{englishcomment}
\begin{chinesecomment}
稿件应明确说明每个神经网络所用输入和输出变量的数量及其物理含义。
\end{chinesecomment}

\Needspace{8\baselineskip}
\subsection{R1.8: Comment 8 / 意见 8}
\begin{englishcomment}
The incorporation of the VIWS formulation into the deep learning framework is not sufficiently explained. It is important to clearly identify where the VIWS formulation enters the neural-network training process and how it contributes to the loss function or governing constraints.
\end{englishcomment}
\begin{chinesecomment}
稿件没有充分解释如何将 VIWS 形式纳入深度学习框架。需要清楚指出 VIWS 形式在神经网络训练过程中的具体作用位置，以及它如何构成损失函数或控制约束。
\end{chinesecomment}

\Needspace{8\baselineskip}
\subsection{R1.9: Comment 9 / 意见 9}
\begin{englishcomment}
It is not entirely clear whether the neutron diffusion equation is enforced through integral residuals derived from the VIWS formulation or through an alternative implementation. This aspect should be clarified.
\end{englishcomment}
\begin{chinesecomment}
目前尚不完全清楚中子扩散方程是通过 VIWS 形式导出的积分残差来施加约束，还是通过其他实现方式施加约束。应澄清这一点。
\end{chinesecomment}

\Needspace{8\baselineskip}
\subsection{R1.10: Comment 10 / 意见 10}
\begin{englishcomment}
Statements such as ``antiderivative forms are used in the deep learning solution'' are difficult for readers to interpret. The manuscript would benefit from a more explicit explanation of how the antiderivative formulation is incorporated into the training procedure and how it contributes to computational efficiency.
\end{englishcomment}
\begin{chinesecomment}
“在深度学习求解中使用原函数形式”之类的表述使读者难以理解。稿件应更明确地解释如何把原函数形式纳入训练过程，以及该形式如何影响计算效率。
\end{chinesecomment}

\Needspace{8\baselineskip}
\subsection{R1.11: Comment 11 / 意见 11}
\begin{englishcomment}
I think the paper is not proposing a new neural network architecture. Rather, it is proposing a new mathematical formulation (VIWS) that is embedded into a deep-learning framework to better handle material interfaces in heterogeneous neutron diffusion problems. If so, please include this explicitly in your manuscript.
\end{englishcomment}
\begin{chinesecomment}
我认为本文并未提出新的神经网络架构，而是提出了一种嵌入深度学习框架的新数学形式（VIWS），以更好地处理非均匀中子扩散问题中的材料界面。如果确实如此，请在稿件中明确说明。
\end{chinesecomment}

\Needspace{8\baselineskip}
\subsection{R1.12: Comment 12 / 意见 12}
\begin{englishcomment}
How do the execution time and memory compare against conventional NDE solvers?
\end{englishcomment}
\begin{chinesecomment}
与传统中子扩散方程（NDE）求解器相比，该方法的执行时间和内存占用如何？
\end{chinesecomment}

\Needspace{14\baselineskip}
\section{Reviewer \#2 / 审稿人 \#2}
\Needspace{8\baselineskip}
\subsection{R2.G1: General Comment 1 / 总体意见 1}
\begin{englishcomment}
The title of the manuscript is a bit misleading. What exactly is ``heterogeneous neutron diffusion calculations''? If the method requires any cross-material homogenization (for example, each pin cell is homogenized, so there are no separate cladding, air gap, fuel and moderator cells), then the word ``heterogeneous'' is not applicable to the manuscript and should be changed to a more appropriate term.
\end{englishcomment}
\begin{chinesecomment}
稿件标题略有误导。“非均匀中子扩散计算”究竟指什么？如果该方法需要进行跨材料均匀化（例如，每个棒元被均匀化，不再分别表示包壳、气隙、燃料和慢化剂区域），那么“非均匀”一词并不适用于本文，应改为更恰当的术语。
\end{chinesecomment}

\Needspace{8\baselineskip}
\subsection{R2.G2: General Comment 2 / 总体意见 2}
\begin{englishcomment}
The quality of the abstract is much lower than the quality of the main body of the manuscript. While the manuscript itself is concise, uses correct terminology and is overall understandable, the abstract does a poor job representing the manuscript. Some sentences in the abstract make no sense at all, so the authors are required to completely rewrite the abstract and carefully proofread it to match the quality of the manuscript's main body.
\end{englishcomment}
\begin{chinesecomment}
摘要的质量明显低于稿件正文。正文较为简洁，术语使用正确，整体上可以理解，但摘要未能很好地反映论文内容。摘要中的某些句子甚至难以理解，因此作者需要完全重写摘要并认真校对，使其质量与正文相匹配。
\end{chinesecomment}

\Needspace{8\baselineskip}
\subsection{R2.G3: General Comment 3 / 总体意见 3}
\begin{englishcomment}
An important aspect of introducing PINNs in place of traditional neutron diffusion solvers is their potential advantage in computational speed. How does the proposed method compare -- apples-to-apples (using exactly the same number of CPU cores/GPU cards, the same models) -- to conventional solvers, and if it is faster, what are the accuracy trade-offs that we can expect? What about memory consumption?
\end{englishcomment}
\begin{chinesecomment}
用 PINN 替代传统中子扩散求解器的一个重要方面，是其潜在的计算速度优势。所提出的方法与传统求解器进行严格同条件比较（使用完全相同数量的 CPU 核/GPU 卡和相同模型）时表现如何？如果更快，需要付出怎样的精度代价？内存消耗又如何？
\end{chinesecomment}

\Needspace{8\baselineskip}
\subsection{R2.G4: General Comment 4 / 总体意见 4}
\begin{englishcomment}
While the presented method seems to be of high value, the potential disadvantages of the method were not clearly stated, so the manuscript reads more like an advertisement, whereas a scientific paper should present both positives and negatives of the work that were observed.
\end{englishcomment}
\begin{chinesecomment}
尽管所提出的方法似乎很有价值，但稿件没有清楚说明其潜在缺点，因此读起来更像宣传材料；科学论文应同时呈现研究中观察到的优点与不足。
\end{chinesecomment}

\Needspace{8\baselineskip}
\subsection{R2.G5: General Comment 5 / 总体意见 5}
\begin{englishcomment}
The proposed method could be introduced in a more reproducible way, which was not the case in the submitted manuscript. The reader will need to dig through all the literature sources mentioned and beyond to piece together the exact pathway on how to implement a similar solver. The manuscript would greatly benefit from having a more straightforward algorithm that the authors used to create the model, train it, and then use in practice. The authors are kindly reminded that they submit their work to the journal named ``Nuclear Engineering and Technology'', and while ``Nuclear'' is clearly present in the manuscript, there is no ``Engineering'' and/or ``Technology'' presented. This makes it look more like a scientific report than an original article. The authors are requested to close this gap.
\end{englishcomment}
\begin{chinesecomment}
所提出的方法应以更具可复现性的方式介绍，而投稿稿件没有做到这一点。读者需要查阅文中提到的全部文献乃至更多资料，才能拼凑出实现类似求解器的具体流程。如果稿件能给出一个更直接的算法，说明作者如何构建模型、训练模型并在实际中使用，将会有很大帮助。还请作者注意，投稿期刊名为《Nuclear Engineering and Technology》。虽然稿件中的“Nuclear”很明确，但没有体现“Engineering”和/或“Technology”，使文章更像一份科学报告，而不是原创研究论文。请作者弥补这一差距。
\end{chinesecomment}

\Needspace{8\baselineskip}
\subsection{R2.S1: Specific Comment 1: Abstract / 具体意见 1：摘要}
\begin{englishcomment}
``Neutron diffusion calculations are essential in reactor physics, but heterogeneous reactor cores contain discontinuous diffusion coefficients and macroscopic cross sections.''

\begin{enumerate}[label=\alph*.,leftmargin=2em]
\item It looks like the authors introduced new and unnecessary terminology in a long-established field. What exactly do the authors mean by ``discontinuous'' diffusion coefficients and ``discontinuous'' macroscopic cross-sections? From the reviewer's view, such combination of words makes no sense.
\item The reviewer is also not convinced that reactor cores contain these coefficients or any other coefficients. The authors are required to proofread the entire Abstract to convey their points in a clearer, more accurate way.
\end{enumerate}
\end{englishcomment}
\begin{chinesecomment}
“中子扩散计算在反应堆物理中至关重要，但非均匀反应堆堆芯包含不连续的扩散系数和宏观截面。”

\begin{enumerate}[label=\alph*.,leftmargin=2em]
\item 看起来作者在一个发展已久的领域中引入了新的且不必要的术语。作者所谓“不连续”的扩散系数和“不连续”的宏观截面究竟是什么意思？在审稿人看来，这样的措辞组合没有意义。
\item 审稿人也不认同“反应堆堆芯包含这些系数或其他系数”的说法。作者需要校对整个摘要，以更加清楚、准确地表达观点。
\end{enumerate}
\end{chinesecomment}

\Needspace{8\baselineskip}
\subsection{R2.S2: Specific Comment 2: Fig. 4 / 具体意见 2：图 4}
\begin{englishcomment}
What is antiderivative, and why are antiderivative fields in Fig. 4 not symmetric?
\end{englishcomment}
\begin{chinesecomment}
什么是原函数？为什么图 4 中的原函数场不对称？
\end{chinesecomment}

\Needspace{8\baselineskip}
\subsection{R2.S3: Specific Comment 3: Section 4.1 / 具体意见 3：第 4.1 节}
\begin{englishcomment}
``The total number of interior and boundary sampling points is 10, 000.''

What are the sampling points and how exactly are they used? How is the number determined and why it matters?
\end{englishcomment}
\begin{chinesecomment}
“域内采样点和边界采样点的总数为 10,000。”

这些采样点是什么？具体如何使用？其数量是如何确定的，为什么这一数量很重要？
\end{chinesecomment}

\Needspace{8\baselineskip}
\subsection{R2.S4: Specific Comment 4: Section 4.3 / 具体意见 4：第 4.3 节}
\begin{englishcomment}
It would be useful to see axially normalized maps and their comparison. In particular, not only the flux map but also the source map that is calculated by the reference code, PINN, and VIWS.
\end{englishcomment}
\begin{chinesecomment}
如果能给出轴向归一化图及其对比，将很有帮助。尤其应同时展示参考程序、PINN 和 VIWS 计算得到的通量图与源项图，而不只是通量图。
\end{chinesecomment}

\Needspace{8\baselineskip}
\subsection{R2.S5: Specific Comment 5: Section 4.4 / 具体意见 5：第 4.4 节}
\begin{englishcomment}
What are the limitations of the method in terms of problem size and boundary condition? Can the meshes be of the entire fuel assembly size?
\end{englishcomment}
\begin{chinesecomment}
该方法在问题规模和边界条件方面有哪些限制？网格能否扩展到整个燃料组件的尺度？
\end{chinesecomment}

\Needspace{8\baselineskip}
\subsection{R2.S6: Specific Comment 6: Section 4.5 / 具体意见 6：第 4.5 节}
\begin{englishcomment}
``The width of each assembly is 20, cm, and the thickness of the surrounding reflector 290 is also 20, cm.''

\begin{enumerate}[label=\alph*.,leftmargin=2em]
\item Commas before the ``cm'' seem unnecessary.
\item Did the VIMS model output pointwise flux shape, or was there a 2D function per each assembly that could produce the flux shape for each point? Please elaborate on the process of flux output in more detail as well as the memory requirement to produce variable resolution flux.
\item Were the fuel assemblies divided into sub-regions? Also, were the fuel assembly regions homogenized material-wise?
\end{enumerate}
\end{englishcomment}
\begin{chinesecomment}
“每个组件的宽度为 20, cm，外围反射层 290 的厚度也是 20, cm。”

\begin{enumerate}[label=\alph*.,leftmargin=2em]
\item “cm”前面的逗号似乎没有必要。
\item VIMS 模型输出的是逐点通量形状，还是每个组件对应一个能够生成各点通量形状的二维函数？请更详细地说明通量输出过程，以及生成不同分辨率通量所需的内存。
\item 燃料组件是否被划分为子区域？此外，燃料组件区域是否在材料参数上进行了均匀化？
\end{enumerate}
\end{chinesecomment}

\Needspace{8\baselineskip}
\subsection{R2.S7: Specific Comment 7: Section 4.5 / 具体意见 7：第 4.5 节}
\begin{englishcomment}
``The effective multiplication factor $k_{\mathrm{eff}}$ is treated as a global parameter to be optimized and is updated together with the neural network weights [45].''

It would be useful to see the details on how the flux and the multiplication factor were produced within the same VIWS model.
\end{englishcomment}
\begin{chinesecomment}
“有效增殖因子 $k_{\mathrm{eff}}$ 被视为待优化的全局参数，并与神经网络权重一起更新 [45]。”

建议说明如何在同一个 VIWS 模型中同时得到通量和增殖因子的具体细节。
\end{chinesecomment}

\end{document}
